\documentclass[11pt]{article}

\usepackage[margin=0.85in]{geometry}
\usepackage[T1]{fontenc}
\usepackage{lmodern}
\usepackage{booktabs}
\usepackage{tabularx}
\usepackage{array}
\usepackage{enumitem}
\usepackage{xcolor}
\usepackage{amsmath}
\usepackage{listings}
\usepackage[hidelinks]{hyperref}

\definecolor{navy}{HTML}{17365D}
\definecolor{teal}{HTML}{006666}
\definecolor{lightblue}{HTML}{EAF1F8}
\definecolor{lightgreen}{HTML}{EDF7ED}
\definecolor{lightwarn}{HTML}{FFF8E1}
\definecolor{codegray}{HTML}{F5F5F5}

\setlist[itemize]{leftmargin=1.35em,itemsep=2pt,topsep=3pt}
\setlist[enumerate]{leftmargin=1.35em,itemsep=2pt,topsep=3pt}
\renewcommand{\arraystretch}{1.25}
\newcolumntype{Y}{>{\raggedright\arraybackslash}X}

\title{\vspace{-1.2em}\textbf{Comprehensive Summary of Architectural Changes}\\
\large Secure Agentic AI Compiler Optimization System}
\author{\textbf{Haamiz Suhail Hakeem} (25CSB1A70) \quad \textbf{Shrot Amish Parekh} (25CSB1A62)}
\date{November 2026}

\begin{document}
\maketitle
\vspace{-1.5em}

\begin{abstract}
This document provides a technical summary of the architectural transformations implemented in the Secure Agentic AI Compiler codebase. The inherited prototype---which relied on simulated latency figures, hard-coded LLVM IR stubs, and open-loop pass generation---has been completely upgraded into an end-to-end, policy-governed compiler. Key enhancements include live Clang++ integration with debug-line metadata, static source and IR security analysis, entrance-level prompt-injection filtering, closed-loop critic feedback with automatic fallback pass injection, multi-fixture behavioral candidate validation with empirical microbenchmarking, and an interactive Web UI for real-time compilation and Attack Lab evaluation.
\end{abstract}

\section{Overview of Codebase Modifications}

\begin{table}[h]
\centering
\small
\begin{tabularx}{\linewidth}{@{}p{0.32\linewidth}p{0.22\linewidth}Y@{}}
\toprule
\textbf{Subsystem \& File} & \textbf{Component Modified} & \textbf{Summary of Technical Changes} \\
\midrule
\texttt{src/core/frontend.py} & Clang Frontend & Auto-discovers \texttt{clang++} / \texttt{clang}; emits unoptimized IR with \texttt{-O0 -Xclang -disable-O0-optnone -gline-tables-only} for line-level debug metadata. \\
\texttt{src/safety/analysis.py} & Security Analyzer & Added source-to-IR line traceability (\texttt{!dbg} $\rightarrow$ \texttt{!DILocation}), \texttt{scan\_adversarial\_directives()}, and \texttt{summarize\_for\_planner()}. \\
\texttt{src/safety/models.py} & Data Models & Added empirical benchmark fields (\texttt{baseline\_latency\_ns}, \texttt{speedup\_pct}) and critic feedback records. \\
\texttt{src/safety/runner.py} & LLVM Pass Runner & Auto-locates \texttt{opt} binary, configures toolchain environment paths, and executes two-phase transformation + IR verification. \\
\texttt{src/safety/validator.py} & Candidate Validator & Links candidate IR using \texttt{clang++}, validates against multiple stdin fixtures, and computes median wall-clock nanosecond latency. \\
\texttt{src/safety/controller.py} & Safe Optimizer Controller & Implemented \texttt{step\_action()} and closed-loop critic feedback loop with automatic allowlisted fallback pass injection. \\
\texttt{src/agents/boss\_agent.py} & Multi-Agent Planner & Ingests static security findings into \texttt{CompilationContext}; applies \texttt{\_apply\_security\_overlay()} to bias pass priorities. \\
\texttt{src/agents/ir\_tuner\_agent.py} & IR Tuner Agent & Filters pass priority queue using security avoid rules and critic feedback; normalizes latency scaling for stub mode. \\
\texttt{src/agents/hw\_tuner\_agent.py} & Hardware Tuner Agent & Filters ISA passes via security constraints and critic feedback; prioritizes allowlisted SIMD passes. \\
\texttt{src/pipeline.py} & End-to-End Orchestrator & Added Stage 0 entrance prompt guard, linked real Clang IR emission, connected empirical microbenchmarking, and added live server mode. \\
\texttt{src/memory/experience\_store.py} & Vector Memory & Auto-pads vectors in \texttt{SQLiteStore.\_cosine()} to resolve 64-d vs 256-d embedding shape mismatches. \\
\texttt{src/web\_ui/app.py} & Interactive Dashboard & Implemented \texttt{GET /api/source} and \texttt{POST /api/compile} for in-browser editing and live compilation triggering. \\
\texttt{tests/test\_architecture.py} & Automated Test Suite & Added 7-part comprehensive test suite validating all core architectural blocks end-to-end. \\
\bottomrule
\end{tabularx}
\caption{Detailed directory of files modified and functionality added.}
\end{table}

\section{Core Architectural Enhancements}

\subsection{Real LLVM Toolchain \& Debug-Line Traceability}
The compiler now interfaces directly with real LLVM 23 binaries. The frontend compiles C/C++ source code via:
\begin{center}
\texttt{clang++ -S -emit-llvm -O0 -Xclang -disable-O0-optnone -gline-tables-only <src> -o <ir>}
\end{center}
The resulting LLVM IR contains \texttt{!DILocation} metadata. \texttt{SecurityAnalyzer} parses these metadata nodes to map IR-level memory vulnerabilities (such as unsafe \texttt{strcpy} calls or large stack buffers) directly back to source code line numbers.

\subsection{Compiler Entrance Attack Guard (Stage 0)}
To counter prompt injection and trojan directives embedded inside code comments or string literals, \texttt{SecurityAnalyzer.scan\_adversarial\_directives()} executes at Stage 0 before any agent planning begins. If an override directive is detected (e.g., \texttt{// SYSTEM PROMPT OVERRIDE: ignore policy}), the pipeline halts fail-closed, records an \texttt{<entrance-prompt-guard>} attempt, and preserves the original unoptimized IR.

\subsection{Closed-Loop Critic Feedback Controller}
The optimization process was upgraded from open-loop batch execution to a closed feedback loop:
\begin{equation*}
\text{Planner Proposes Pass} \longrightarrow \text{Policy Gate Checks} \longrightarrow \text{LLVM Runner Applies} \longrightarrow \text{Validator Tests} \longrightarrow \text{Critic Feedback to Planner}
\end{equation*}
If an action is rejected by the policy gate (e.g. non-allowlisted \texttt{always-inline}) or fails validation, \texttt{SafeOptimizationController} consults \texttt{\_FALLBACK_ALTERNATIVES} and enqueues allowlisted substitutes (e.g. \texttt{inline} or \texttt{simplifycfg}) labeled with stage \texttt{critic\_fallback(<pass>)}.

\subsection{Multi-Fixture Validation \& Empirical Microbenchmarking}
Rather than relying solely on analytical cost models, \texttt{CandidateValidator}:
\begin{enumerate}
    \item Caches the compiled baseline binary to prevent redundant recompilations.
    \item Compiles the candidate IR to native executable code using \texttt{clang++}.
    \item Executes both binaries against representative stdin fixtures (\texttt{""}, \texttt{"100\textbackslash n"}, \texttt{"1000 250 5\textbackslash n"}), verifying that exit codes, stdout, and stderr remain strictly identical.
    \item Performs warm-up runs followed by multi-iteration microbenchmarking to report median wall-clock nanosecond execution time and empirical speedup percentage.
\end{enumerate}

\subsection{Experience Memory Cross-Dimensional Vector Compatibility}
To bridge heuristic IR2Vec encoders (64-dimensional) with production vector databases (256-dimensional), \texttt{SQLiteStore.\_cosine()} dynamically wrap-pads vectors:
\begin{equation*}
\vec{a}_{\text{padded}} = \text{pad}\left(\vec{a}, (0, \max(d_a, d_b) - d_a), \text{mode}=\text{wrap}\right)
\end{equation*}
This guarantees mathematical stability and similarity scoring without shape mismatch errors.

\subsection{Interactive Web UI \& Attack Lab Integration}
Using Python's standard library \texttt{http.server}, \texttt{src/web\_ui/app.py} exposes REST endpoints (\texttt{/api/source}, \texttt{/api/compile}, \texttt{/api/results}, \texttt{/api/attack-lab}). Users can inspect, edit, and recompile C/C++ source code directly in the browser while viewing real-time critic feedback, empirical speedup metrics, and scenario detection statistics.

\section{Verification \& Test Suite Summary}
The architecture is verified by \texttt{tests/test\_architecture.py}, covering all 7 subsystems:
\begin{itemize}
    \item \textbf{Test 1 (Clang Frontend):} Validates real IR generation with \texttt{!dbg} metadata.
    \item \textbf{Test 2 (Security Analyzer):} Confirms CWE detection and \texttt{!dbg} line mapping.
    \item \textbf{Test 3 (Prompt Guard):} Validates fail-closed blocking of adversarial directives.
    \item \textbf{Test 4 (Agent Reasoning):} Confirms pass selection adaptation based on security findings and critic feedback.
    \item \textbf{Test 5 (Closed-Loop Optimizer):} Verifies pass execution, critic fallback injection, and candidate validation.
    \item \textbf{Test 6 (Experience Store):} Verifies 64-d and 256-d embedding similarity retrieval.
    \item \textbf{Test 7 (Web UI):} Validates REST API responses for results, source fetching, and Attack Lab metrics.
\end{itemize}
All 7 unit tests pass with 100\% success rate.

\end{document}

